% Compilar com o XeLaTeX
\documentclass[a5paper,twoside,openany,10pt]{memoir}

\usepackage{graphicx}
\usepackage{xcolor}
\usepackage{hyperref}
\usepackage{memhfixc}

% Set font to Adobe Palatino
\usepackage{fontspec}
\usepackage{mathpazo}
\setmainfont
     [ BoldFont       = texgyrepagella-bold.otf ,
       ItalicFont     = texgyrepagella-italic.otf ,
       BoldItalicFont = texgyrepagella-bolditalic.otf ]
     {texgyrepagella-regular.otf}

%\renewcommand{\rmdefault}{ppl}
%\renewcommand{\sfdefault}{phv}
%\renewcommand{\ttdefault}{pcr}

\newcommand{\mailto}[1]{\href{mailto:{#1}}{#1}}
\newcommand{\link}[2]{\href{#2}{#1}}

\title{PSI3211 - Circuitos I}
\author{G. Queiroz\\<\mailto{gabrieljvnq@gmail.com}>}
\date{2019}

\begin{document}
\maketitle

$x = \frac{-b \pm \sqrt{b^2 - 4ac}}{a}$

Luís argüia à Júlia que «brações, fé, chá, óxido, pôr, zângão» eram palavras do português.

À noite, vovô Kowalsky vê o ímã cair no pé do pinguim queixoso e vovó põe açúcar no chá de tâmaras do jabuti feliz. (90 letras, incluindo todas as letras acentuadas)

\begin{center}
\includegraphics[width=1cm]{CC-BY-SA.ps}\\
\footnotesize
Documento sob a \link{Licença Creative Commons Atribuição-CompartilharIgual}{https://creativecommons.org/licenses/by-sa/4.0/}.
\end{center}

\end{document}
